\section{A positive charge interval}
\label{subsec:interacting-positive-threshold}

We now verify that the response conditions admit
a common interval of charges, independent of the
ultraviolet cutoff. In addition to
\eqref{eq:interacting-factor-conditions},
\eqref{eq:interacting-small-complex-mass},
\eqref{eq:interacting-one-cell-conditions},
\eqref{eq:interacting-tree-smallness}, and
\eqref{eq:interacting-jet-conditions}, require
$r_*>31L+16$, $v_*\ge8$, $p_0v_*^{p_0-1}\ge28$, and
\begin{gather}
 Mr_*/(8L)-1\ge6M+8,\qquad
 \lambda_L\le\bar\kappa_1^s,
                                  \label{eq:interacting-last-response-conditions}\\
 D_\beta^{\rm buf}/2+1/2\le3\ell_c/8,\qquad
 D_\beta^{\rm buf}+3\le\ell_c .
                                  \label{eq:interacting-placement-conditions}
\end{gather}
The construction of \cite[\S6.1]{CoverTuning}
fixes an odd $L\ge13$, a power
$M$ of $L$ with $M\ge L^2$, a localization
power $r_{\rm loc}\ge2$, and cutoff powers
$p_0\ge2$, $p>2p_0$. It also fixes
$0<\varepsilon\le1/200$ and $0<s<1$ such that
\begin{equation}\label{eq:interacting-response-exponents}
 \mathfrak e=s(1/4-9\varepsilon)-12\varepsilon>0.
\end{equation}
For the construction fixed in the introduction, take $s=1/2$;
then $\mathfrak e\ge17/400$.

\begin{proposition}\label{prop:interacting-positive-threshold}
For the fixed construction constants, $\bar m$,
and probe $\beta$, there is $e_\sharp(\beta)>0$
such that all response conditions hold for
$0<e\le e_\sharp(\beta)$ and every
$N\ge N_{\rm ref}(e)$. The threshold depends
on the probe only through
$C_\beta,\mathbf D,D_\beta^{\rm buf}$.
It is independent of $N$, the twist, and the
admitted placement. All centered placements
are admitted in this range.
\end{proposition}

\begin{proof}
Write
$v_K=-\log(e\sqrt h)$ and
$r_K=\min\{L^i:L^i\ge v_K^{r_{\rm loc}}\}$.
Then $R^*(e)=hMr_K/L$ and
$\ell_c=R^*(e)/(2L)$.
\cite[Proposition~6.1]{CoverTuning} and
\cite[(53)]{CoverTuning} give
a finite $N_{\rm ref}(e)$.

For $N\ge N_{\rm ref}(e)$ the stop level
is $K-1$ or $K-2$ by
\cite[(54)]{CoverTuning}. In particular its
scales are independent of $N$ and satisfy
\begin{gather}
 v_K\le v_*,\quad e_*=\mathrm e^{-v_*},\quad
 a_*\in\{h/L,h/L^2\},\notag\\
 v_*^{r_{\rm loc}}\le r_*,\quad
 Mv_*^{r_{\rm loc}}\le\ell_*<LMv_*^{r_{\rm loc}},
 \quad p_{0*}=v_*^{p_0},\quad p_*=v_*^p,\notag\\
 p_\beta\le p_\beta^+:=
                (D_\beta^{\rm buf}L^2/h+2M+4)^3,\qquad
 \ell_c\ge hMv_K^{r_{\rm loc}}/(2L^2).
                                  \label{eq:interacting-threshold-scales}
\end{gather}
The strict wrapping condition $r_*>31L+16$
already follows from \cite[(50)]{CoverTuning}
whenever $e\sqrt h\le e_1$.
Thus every response condition is a condition
on $e$ alone once the construction and probe
are fixed.

A bound $A v^d \mathrm e^{-\gamma v}$
with $A,d\ge0$ and $\gamma>0$ tends to zero,
uniformly on a sufficiently large half-line.
Split the exponential in half and use
$v^d\mathrm e^{-\gamma v/2}\le(2d/(\gamma\exp(1)))^d$.
Likewise, $cv^2$ exceeds $bv+\log^+A$
on a sufficiently large half-line for
$c>0$. We apply these two inequalities to
the following groups of conditions.

The local tree smallness and terminal factor
conditions follow from exponential decay
with positive rates
$\varepsilon$, $2\varepsilon$,
$1/4-\varepsilon$, and $1/4-9\varepsilon$;
the volume factors are bounded by
$(LM)^3v_*^{3r_{\rm loc}}$.
For $\bar r$, use
$1/4-7\varepsilon>2\varepsilon$.
The complex mass and coarse determinant
conditions similarly have rate $2\varepsilon$.
The support count is bounded by $p_\beta^+$.
All three summands of $B_0$ tend to zero:
their rates are $4\varepsilon$, $2\varepsilon$,
and the quadratic decay $\mathrm e^{-r_*}\le \mathrm e^{-v_*^2}$.
Also $\tau_s\ge c_\tau Mv_*^2$, so both
$\tau_s\ge1$ and
$52B_1\mathrm e^{-\tau_s}\le e_*^{1/4-9\varepsilon}$
hold once $v_*$ is large.

The separation condition follows from
$r_*\ge v_*^2$. Its error obeys
\[
 2\varpi_{\rm sep}^+
 \le2^{N_*+1}C_W^{\mathrm{mix}}p_\beta^+
          \mathrm e^{\vartheta_0(4M+7)}
          \mathrm e^{-\vartheta_0Mv_*^2/(8L)} .
\]
Since $\bar\kappa_1\ge \mathrm e^{-v_*/4}$,
quadratic decay makes this at most
$\bar\kappa_1^{N_*}$ for all sufficiently
large $v_*$. The family summation condition
$N_*\bar\kappa_1^{1-s}\le1/2$
uses the positive rate
$(1-s)(1/4-9\varepsilon)$.

For the photon complement, the inequalities $v_*\ge8$ and
$p_0v_*^{p_0-1}\ge28$ hold on a sufficiently large half-line,
since $p_0\ge2$. Moreover $p>2p_0$ implies
$(p_*-1)^2\ge C_0^2p_{0*}$ once $v_*$
is large. Therefore
\begin{equation}\label{eq:interacting-complement-threshold}
 \varpi_0\le24(LM)^3v_*^{3r_{\rm loc}}\mathrm e^{-v_*^2/2}.
\end{equation}
This is eventually at most $1/4$.
Moreover $t_2/e$ is fixed and positive, while
$e$ is between $(L/h)^{1/2}e_*$
and $Lh^{-1/2}e_*$. Hence
$8N_*t_2/t_x$ is bounded by a fixed
constant times $e_*^{1-3\varepsilon}$,
which tends to zero.

The potentially largest loss is the fourth
derivative radius in $\lambda_L$.
From \eqref{eq:interacting-complement-threshold},
the definition of $\lambda_L$, and
$r/t_2\le C \mathrm e^{v_*}$, we obtain
\begin{equation}\label{eq:interacting-lambda-threshold}
 \lambda_L\le A_8
   \exp\left(-\frac{v_*^2}{12N_*}
                     +\left(\frac{r_{\rm loc}}{2N_*}+4\right)v_*\right)
\end{equation}
for a finite $A_8>0$ depending only on the
fixed constants and probe bounds.
For clarity one can take
\[
 A_8=[24(LM)^3]^{1/(6N_*)}
       \mathrm e^{\mathbf D/(6N_*)}M_L^{2/(3N_*)}
       5\left[
        \frac{N_*\sqrt h}{2C_b\sqrt L}(t_2/e)^{-1}\right]^4 .
\]
If
$v_*^2/(12N_*)\ge
(r_{\rm loc}/(2N_*)+17/4)v_*+\log^+A_8$,
then $\lambda_L\le \mathrm e^{-v_*/4}\le\bar\kappa_1^s$.
Thus the complex continuation and the
fourth derivatives preserve a positive
small-charge interval.
Finally the two placement inequalities
hold when
$\ell_c\ge(4/3)D_\beta^{\rm buf}+3$;
this follows from
\eqref{eq:interacting-threshold-scales} as $v_K\to\infty$.

Choose a finite $V$ beyond all the half-lines
just specified and put
$e_V=h^{-1/2}\min(e_1,\mathrm e^{-V})>0$.
Every $0<e\le e_V$ has $v_*\ge v_K\ge V$
and satisfies every response condition.
Define $e_\sharp$ as one half of the
supremum of the charges $e'\le e_1h^{-1/2}$
for which all conditions hold for every
$0<e\le e'$. Its defining set contains
$e_V$, so $e_\sharp\ge e_V/2>0$.
For $e\le e_\sharp$, a member of the
defining set exceeds $e_\sharp$, which
ensures that all inequalities hold at $e$
even if the scale functions have jumps.
This avoids assuming monotonicity of the
rounded stop scales.

The scale identities depend on $e,h$ and
the fixed construction constants; the probe
enters through
$C_\beta,\mathbf D,D_\beta^{\rm buf}$.
They do not involve the twist or placement,
and they are unchanged by the coordinate
reflections used to form the cell probes.
Centered placements satisfy the admission
inequality by
\eqref{eq:interacting-placement-conditions},
since their centers differ from the cell
center by at most $\eta/2\le1/2$.
The same conditions give
$D_\beta^{\rm buf}+3\eta\le\ell_c$ and
$R^*(e)\ge D_\beta+1$.
\end{proof}



\begin{proof}[Proof of \cref{thm:rrm-main}]
The base construction gives the positive physical ceiling
\[
 e_G=\begin{cases}
 e_1,&h=1,\\
 \min\{e_1/\sqrt h,\,1/(\sqrt h\log h)\},&h>1,
 \end{cases}
\]
as in \cite[(31)]{QED3Integration}. On this interval the common pair
and all upstream terminal estimates apply.
Choose $e_{\rm RR}=\min\{e_G,e_\sharp(\beta)\}$, with the threshold
of \cref{prop:interacting-positive-threshold}, at $s=1/2$.
That proposition verifies the additional inequalities of
\cref{thm:interacting-reference-error} for every smaller positive
charge and every $N\ge N_{\rm ref}(e)$, and admits every centered
placement. The bounds \eqref{eq:interacting-final-terminal-error}
give \eqref{eq:rrm-main-bound} with the stated common constant.
The exact unfolding identifies the normalized density with the
zero-flux law, \cref{prop:external-entire} proves entireness, and
\eqref{eq:external-Wick-product} proves \eqref{eq:rrm-main-wick}.
Its differentiation proof establishes the required integrability.
\end{proof}
